\documentclass[a4paper,11pt]{article} % Don't change this
\usepackage{fullpage} % Don't change this

\usepackage{titling} % Don't change this
\pretitle{\vspace{-6em}\begin{center}\larger[2]\bf} % Don't change this
\postauthor{\end{center}} % Don't change this
\preauthor{\vspace{-2em}\begin{center}\smaller[1]} % Don't change this
\postauthor{\end{center}\vspace{-5em}} % Don't change this

\usepackage[shortlabels]{enumitem} % Don't change this 
\setlist{itemsep=0em,topsep=0.3em} % Don't change this

%%% Some useful packages.
\usepackage{xcolor}
\usepackage{float}
\usepackage{amsmath,amssymb,mathtools}
\usepackage{graphicx}
\usepackage{subcaption}
\usepackage[hidelinks]{hyperref}
\usepackage{relsize}
\usepackage{cleveref}
\crefname{figure}{Fig.}{Figs.}

\usepackage[backend=bibtex,natbib=true,style=authoryear-icomp,maxnames=3,giveninits=true]{biblatex}
\citetrackerfalse
\bibliography{ref} % put bib entries in ref.bib

\usepackage[linesnumbered]{algorithm2e}
\usepackage{cleveref}
\crefname{equation}{Equation.}{Equations.}
\crefname{figure}{Figure}{Figures.}

\usepackage{lipsum}
\usepackage{wrapfig}
\usepackage{tikz}

%%% Add additional packages or macros below

\begin{document}

\title{STAT3007/7007 Project Proposal Template}
\author{Group X \\
Issac Newton (s12345678) $\cdot$
Albert Einstein (s23456781) $\cdot$
Alan Turing (s34567812) \\
John von Neumann (s45678123) $\cdot$ Kurt G\"odel (s56781234)}

\date{}
\maketitle

\begin{abstract}
	This is a \LaTeX\ template for your project proposal.
	Note the following: 
	(a) Change the title and group information above to your project title and
	group information.
	(b) Your proposal should not exceed \underline{2 pages} (excluding the
	references) 
	You should not alter the font size or the page size.
	(c) You can do collaborative writing by uploading the template to \underline{Overleaf}:
	download \emph{STAT3007-7007\_Project\_Templates.zip},
	login to \url{https://www.overleaf.com},
	click \emph{Upload Project},
	then upload the zip file.
	Overleaf provides excellent guides on \LaTeX\ at
	\url{https://www.overleaf.com/learn}.
	Often, a quicker way to get the answer to a \LaTeX\ question is to ask ChatGPT;
	in many cases, it can give you a correct answer without much prompting.
	(d) When you submit your proposal, name it in the format of
	\emph{GX\_proposal\_AlanTuring.pdf}, where 
	\emph{GX} should be replaced by your group name, and \emph{AlanTuring} by the
	submitter's name.
	Only one group member needs to submit.
\end{abstract}

\section{Topic}

Briefly describe your project topic.
This can be as short as one or two sentences, but you should make it clear what
the machine learning problem is.

\section{Significance}

Briefly explain why the problem is important.
See \emph{report.tex} for how to cite references.

\section{Feasibility}

This is the most substantial part of your proposal.
You should justify that the project is feasible by explaining what has been
done, what the planned work will be and how it will be carried out.  
If data is needed for your project, make sure that you have found some, as
sometimes data does not come as easily as we expect.

\smaller
\printbibliography

\end{document}
